% Verbatim from minor_report/chapters/chapter_3_methodology.tex lines 1142-1211 (retired report),
% headings demoted by 0 level(s). Do not edit here; edit the source of the numbers.

VINS-Fusion uses stereo features and inertial measurements in a sliding-window
nonlinear estimator. The ROS wrapper synchronizes the raw image topics and uses
separate callback groups for images, IMU, feature messages, and detections. The
simulation configuration processes every image (\repo{image_skip: 1}), uses two
cameras and IMU, fixes simulated extrinsics, targets \SI{25}{\hertz} feature
publication, and uses 10-pixel keyframe parallax. Image queues were increased to
100 and the IMU queue to 2,000 to reduce losses during initialization and backlog.

VINS cannot use only the ORB strategy because it propagates existing features
with pyramidal KLT optical flow and detects new features only for replenishment.
A feature attached to an object before detection could otherwise persist after
masking. The implementation initializes \texttt{FeatureTracker::setMask()} from
the YOLO mask. Tracks in zero-valued regions fail the retention test, while the
same mask is passed to \texttt{goodFeaturesToTrack()} to suppress new points. The
count of persistent tracks removed in dynamic regions is accumulated as a
mechanism diagnostic.

\begin{figure}[htbp]
  \centering
  \resizebox{\textwidth}{!}{% VINS-Fusion pipeline, drawn to show that the mask must act at TWO points.
% The contrast with the ORB figure is the point of both: ORB extracts fresh
% features every frame, so rejecting candidates at extraction is sufficient.
% VINS propagates existing KLT tracks, so suppressing new features alone would
% leave a track that latched onto an object before it was detected still
% attached to it. Both insertions are therefore required.
%
% Layout note: marginalization sits ABOVE the optimizer rather than below it so
% that the IMU can be routed orthogonally beneath the whole diagram. A straight
% IMU-to-optimizer line cuts through goodFeaturesToTrack() and the highlight box.
\begin{tikzpicture}[
  font=\footnotesize,
  box/.style={draw, rounded corners, align=center, minimum height=8mm,
    text width=26mm, inner sep=1.6mm, fill=implemented},
  up/.style={box, fill=gray!12},
  det/.style={box, fill=external},
  ext/.style={box, fill=external},
  io/.style={box, fill=gray!5, draw=gray!60},
  flow/.style={-{Latex[length=2mm]}, thick},
  opt/.style={flow, dashed},
  lbl/.style={font=\scriptsize, inner sep=1pt}
]
  \node[det] at ( 0.0,  4.6) (yolo)   {YOLO detector\\(optional)};
  \node[box, text width=32mm]
             at ( 4.8,  4.6) (setmask){\texttt{FeatureTracker::}\\\texttt{setMask()}};
  \node[up]  at ( 4.8,  2.3) (klt)    {propagate tracks\\pyramidal KLT};
  \node[up]  at ( 4.8,  0.0) (gftt)   {\texttt{goodFeatures}\\\texttt{ToTrack()}};
  \node[io]  at ( 0.0,  1.15)(stereo) {stereo pair\\\SI{30}{\hertz}};
  \node[io]  at ( 0.0, -2.6) (imu)    {IMU\\\SI{200}{\hertz}};
  \node[up]  at ( 9.8,  1.15)(window) {sliding window\\optimization};
  \node[up]  at ( 9.8,  3.8) (marg)   {marginalization};
  \node[io]  at (14.4,  1.15)(out)    {\repo{vio.csv},\\\texttt{world}$\to$\texttt{body}};
  \node[ext] at (14.4, -2.6) (rtab)   {RTAB-Map\\external back end};

  \draw[flow] (stereo.east) -- (klt.west);
  \draw[flow] (stereo.east) -- (gftt.west);
  \draw[flow] (klt.east)    -- (window.west);
  \draw[flow] (gftt.east)   -- (window.west);
  % Orthogonal, beneath everything: nothing sits between y=-2.6 and the
  % optimizer's south face at x=9.8.
  \draw[flow] (imu)         -| (window.south);
  \draw[flow] ([xshift=-4mm]window.north) -- ([xshift=-4mm]marg.south);
  \draw[flow] ([xshift=4mm]marg.south)    -- ([xshift=4mm]window.north);
  \draw[flow] (window)      -- (out);
  \draw[flow] (out)         -- (rtab);
  \draw[opt]  (yolo)        -- (setmask);

  % --- the two insertions, which is the reason the figure exists
  \draw[opt] (setmask) -- (klt);
  \node[lbl, right=1mm] at (4.8, 3.45) {\textbf{1} cull};
  % Down the corridor between the highlight box and the optimizer.
  \draw[opt] (setmask.east) -- (7.3, 4.6) -- (7.3, 0.0) -- (gftt.east);
  \node[lbl, right=1mm] at (7.3, 2.6) {\textbf{2} suppress};
  \node[draw=flagred, thick, rounded corners, inner sep=2.4mm,
        fit=(klt)(gftt)] (hl) {};
  \node[lbl, text=flagred, below=1.5mm of hl] {both insertions required};
\end{tikzpicture}
}
  \caption[VINS-Fusion pipeline and its two semantic-mask insertion points.]{VINS-Fusion pipeline and its two semantic-mask insertion points.
  Because the front end propagates existing KLT tracks and only detects new
  features to replenish them, suppressing new detections alone is insufficient:
  a track that attached to an object before it was detected would survive
  masking. Tracks in masked regions must therefore also be culled (1) while the
  same mask suppresses replenishment (2). Contrast \cref{fig:orb-pipeline},
  where re-extraction every frame makes a single insertion sufficient. The count
  of persistent tracks removed is logged as a mechanism diagnostic --- the
  counter that read zero in all twelve filtered runs.}
  \label{fig:vins-pipeline}
\end{figure}

\begin{table}[htbp]
\centering
\caption[VINS-Fusion simulation configuration at the cutoff, from \repo{vins_fusion_ros2/config/wil_sim/stereo_imu.yaml}.]{VINS-Fusion simulation configuration at the cutoff, from
\repo{vins_fusion_ros2/config/wil_sim/stereo_imu.yaml}. Inertial noise densities
are in \cref{tab:imu-noise-assumed}.}
\label{tab:vins-params}
\small
\setlength{\tabcolsep}{4pt}
\begin{tabularx}{\textwidth}{>{\raggedright\arraybackslash}p{0.26\textwidth}l>{\raggedright\arraybackslash}X}
\toprule
Parameter & Value & Role \\
\midrule
\repo{num_of_cam} / \repo{imu} & 2 / 1 & stereo-inertial \\
\repo{image_skip} & 1 & every image processed \\
\repo{max_cnt} & 150 & tracked-feature budget \\
\repo{min_dist} & \SI{30}{px} & minimum feature separation \\
\repo{freq} & \SI{25}{\hertz} & target feature publication \\
\repo{keyframe_parallax} & \SI{10}{px} & keyframe decision \\
\repo{F_threshold} & 1.0 & RANSAC fundamental-matrix gate \\
\repo{max_solver_time} & \SI{0.08}{\second} & Ceres budget per update \\
\repo{max_num_iterations} & 100 & Ceres iteration cap \\
\repo{estimate_extrinsic} & 0 & simulated extrinsics held fixed \\
\repo{estimate_td} & 0 & camera--IMU delay held at \SI{0}{\second} \\
\repo{g_norm} & \SI{9.8}{\metre\per\second\squared} & matches the world's gravity \\
Image queue & 100 & raised to survive initialization backlog \\
IMU queue & \num{2000} & raised for the same reason \\
\bottomrule
\end{tabularx}
\end{table}

VINS writes the optimized state after each estimator update to \repo{vio.csv} and
broadcasts \texttt{world~$\rightarrow$~body} using configurable frame names. The
ROS \repo{output_path} parameter overrides the YAML path, recursively creates
missing directories, and reports an explicit error if \repo{vio.csv} cannot be
opened. This behaviour is an uncommitted worktree change at the cutoff. Evidence
is in \repo{vins_fusion_ros2/src/vins_estimator.cpp},
\repo{vins_fusion_ros2/vins/src/featureTracker/feature_tracker.cpp}, and
\repo{vins_fusion_ros2/vins/src/estimator/estimator.cpp}.
